\kchapter{Observability: Logs, Metrics, Traces}{prometheus}

\begin{objectives}
\begin{itemize}
  \item the three pillars of observability and how they fit together;
  \item how container logging works and how to centralise logs;
  \item metrics with Prometheus and dashboards with Grafana;
  \item alerts that wake you up only when it matters;
  \item distributed tracing with OpenTelemetry.
\end{itemize}
\end{objectives}

\section{Why observability?}

In a cluster, Pods move, scale, and die constantly. You cannot SSH into ``the
server'' to look around. You need the system to \emph{tell you} what is going
on. That ability is \kterm{observability}, built from three signals:

\begin{center}
\resizebox{0.95\linewidth}{!}{%
\begin{tikzpicture}[
  card/.style={draw=#1, fill=#1!8, rounded corners=6pt, line width=1pt,
               minimum width=4.6cm, minimum height=3.1cm, align=center, text width=4.2cm,
               font=\footnotesize, text=k8snavy}]
  \node[card=k8sblue] at (0,0) {{\Large\color{k8sblue}\faStream}\\[3pt]{\large\bfseries Logs}\\[3pt]
     \emph{What happened?}\\Timestamped text events.\\[2pt]\scriptsize Loki, Elasticsearch, CloudWatch};
  \node[card=k8sorange] at (5.2,0) {{\Large\color{k8sorange}\faChartLine}\\[3pt]{\large\bfseries Metrics}\\[3pt]
     \emph{How much / how often?}\\Numbers over time.\\[2pt]\scriptsize Prometheus, Grafana, Mimir};
  \node[card=k8spurple] at (10.4,0) {{\Large\color{k8spurple}\faProjectDiagram}\\[3pt]{\large\bfseries Traces}\\[3pt]
     \emph{Where did the time go?}\\One request across services.\\[2pt]\scriptsize OpenTelemetry, Tempo, Jaeger};
\end{tikzpicture}}
\end{center}

\section{Logs}

Containers write to \kc{stdout} and \kc{stderr}. The container runtime saves
that output to files on the node, and the kubelet serves them to
\kc{kubectl logs}.

\begin{bashcode}{terminal}
kubectl logs deploy/api                     # one Pod of the Deployment
kubectl logs -l app=api --all-containers --prefix --since=15m
kubectl logs api-7d9f8c-xk2lp --previous    # the run before the last crash
kubectl logs -f api-7d9f8c-xk2lp
stern api -n shop                           # tail all matching Pods, colourised
\end{bashcode}

\begin{warningbox}
Node log files are rotated and disappear with the node. For anything beyond
live debugging, ship logs to a central store: a DaemonSet agent
(\textbf{Fluent Bit}, \textbf{Grafana Alloy}, \textbf{Vector}) on each node reads
the log files and forwards them to \textbf{Loki}, \textbf{Elasticsearch/OpenSearch},
or your cloud's logging service.
\end{warningbox}

\begin{protip}
Log in \textbf{structured JSON} with fields like \kc{level}, \kc{msg},
\kc{trace\_id}, and \kc{user\_id}. Structured logs can be filtered and
correlated; free text can only be grepped.
\end{protip}

\section{Metrics with Prometheus}

\kicon{prometheus}~\textbf{Prometheus} is the CNCF standard for metrics. It
\emph{pulls} (``scrapes'') an HTTP endpoint, usually \kc{/metrics}, on every
target every 15--30 seconds and stores the numbers as time series.

\begin{plaincode}{GET /metrics}
# HELP http_requests_total Total HTTP requests.
# TYPE http_requests_total counter
http_requests_total{method="GET",handler="/api/tasks",status="2xx"} 10432
http_requests_total{method="POST",handler="/api/tasks",status="5xx"} 3
# HELP http_request_duration_seconds Request latency.
# TYPE http_request_duration_seconds histogram
http_request_duration_seconds_bucket{handler="/api/tasks",le="0.1"} 9870
\end{plaincode}

\subsection{The kube-prometheus-stack}

The easiest way to get a complete monitoring setup is the
\textbf{kube-prometheus-stack} Helm chart. It installs:

\begin{itemize}
  \item the \textbf{Prometheus Operator} and a Prometheus server;
  \item \textbf{Alertmanager} for routing alerts (Slack, PagerDuty, e-mail);
  \item \textbf{Grafana} with ready-made dashboards for nodes, Pods, and workloads;
  \item \textbf{node-exporter} (node metrics) and \textbf{kube-state-metrics}
        (object state: replicas, Pod phases, restarts).
\end{itemize}

\begin{bashcode}{terminal}
helm repo add prometheus-community https://prometheus-community.github.io/helm-charts
helm repo update
helm install monitoring prometheus-community/kube-prometheus-stack \
  --namespace monitoring --create-namespace

kubectl -n monitoring port-forward svc/monitoring-grafana 3000:80
# open http://localhost:3000  (user: admin; password below)
kubectl -n monitoring get secret monitoring-grafana \
  -o jsonpath='{.data.admin-password}' | base64 -d
\end{bashcode}

\subsection{Telling Prometheus what to scrape}

With the operator, you describe scrape targets with a \kc{ServiceMonitor}:

\begin{yamlcode}{api-servicemonitor.yaml}
apiVersion: monitoring.coreos.com/v1
kind: ServiceMonitor
metadata:
  name: api
  namespace: shop
  labels:
    release: monitoring        # the chart's default selector picks this up
spec:
  selector:
    matchLabels:
      app: api                 # the Service to scrape
  endpoints:
    - port: http               # named Service port
      path: /metrics
      interval: 15s
\end{yamlcode}

\subsection{PromQL in five queries}

\begin{plaincode}{PromQL}
# Requests per second, per handler, over the last 5 minutes
sum by (handler) (rate(http_requests_total[5m]))

# Error ratio (5xx / all)
sum(rate(http_requests_total{status="5xx"}[5m])) / sum(rate(http_requests_total[5m]))

# 95th percentile latency
histogram_quantile(0.95, sum by (le) (rate(http_request_duration_seconds_bucket[5m])))

# Container restarts in the last hour
increase(kube_pod_container_status_restarts_total[1h]) > 0

# Memory usage as % of limit
container_memory_working_set_bytes / on(pod, container) kube_pod_container_resource_limits{resource="memory"}
\end{plaincode}

\section{What to measure: golden signals}

\begin{center}
\ktablestart
\begin{tabularx}{\linewidth}{P{2.4cm} L}
\kheader \kth{Signal} & \kth{Question} \\
\textbf{Latency} & How long do requests take (p50, p95, p99)? \\
\textbf{Traffic} & How many requests per second? \\
\textbf{Errors} & What fraction of requests fail? \\
\textbf{Saturation} & How full is the system (CPU, memory, queue depth, connections)? \\
\end{tabularx}
\end{center}

\section{Alerts}

\begin{yamlcode}{api-alerts.yaml}
apiVersion: monitoring.coreos.com/v1
kind: PrometheusRule
metadata:
  name: api-alerts
  namespace: shop
  labels:
    release: monitoring
spec:
  groups:
    - name: api.rules
      rules:
        - alert: APIHighErrorRate
          expr: |
            sum(rate(http_requests_total{namespace="shop",status="5xx"}[5m]))
              / sum(rate(http_requests_total{namespace="shop"}[5m])) > 0.05
          for: 10m
          labels:
            severity: page
          annotations:
            summary: "API error rate above 5% for 10 minutes"
            runbook_url: "https://wiki.example.com/runbooks/api-errors"
\end{yamlcode}

\begin{protip}
Alert on \textbf{symptoms users feel} (error rate, latency, availability ---
ideally expressed as SLOs), not on every cause (``CPU is 80\%''). Every paging
alert should be actionable and link to a runbook.
\end{protip}

\section{Traces with OpenTelemetry}

A single click may touch a gateway, three services, a cache, and a database. A
\kterm{trace} follows that request end-to-end as a tree of timed
\kterm{spans}. \textbf{OpenTelemetry} (OTel) is the vendor-neutral standard:
instrument once with the OTel SDK, send data to an \textbf{OTel Collector}
running in the cluster, and export to Tempo, Jaeger, or a commercial backend.

\begin{yamlcode}{otel-env.yaml}
env:
  - name: OTEL_SERVICE_NAME
    value: taskflow-api
  - name: OTEL_EXPORTER_OTLP_ENDPOINT
    value: http://otel-collector.observability:4317
  - name: OTEL_RESOURCE_ATTRIBUTES
    value: deployment.environment=production
\end{yamlcode}

\section{Kubernetes events}

Events are the cluster's own log: scheduling decisions, image pulls, probe
failures, OOM kills. They expire after about an hour, so export them
(kube-state-metrics, an event exporter) if you want history.

\begin{bashcode}{terminal}
kubectl get events -n shop --sort-by=.lastTimestamp
kubectl get events -A --field-selector type=Warning
\end{bashcode}

\begin{recap}
\begin{itemize}
  \item Logs: stdout/stderr, shipped by a node agent to a central store.
  \item Metrics: Prometheus scrapes \kc{/metrics}; Grafana visualises; Alertmanager notifies.
  \item ServiceMonitors and PrometheusRules configure monitoring declaratively.
  \item Traces: OpenTelemetry shows where time goes across services.
\end{itemize}
\end{recap}
